% The butterfly network an FFT runs on: at stage s every point combines with
% the point 2^s away, so after log2(N) stages every output depends on every
% input -- which is why the communication pattern is effectively all-to-all.
\documentclass[dvisvgm,border=3pt]{standalone}
\input{preamble.tex}
\begin{document}
\begin{tikzpicture}[x=1cm,y=1cm,
    pt/.style ={circle,fill=uibkorange,minimum size=3.6mm,inner sep=0pt},
    wire/.style={draw=blu,line width=0.6pt}]

  \def\dx{1.35}   % distance between stages
  \def\dy{0.62}   % distance between points

  % four stages of eight points
  \foreach \s in {0,...,3}{
    \foreach \i in {0,...,7}{
      \node[pt] (p\s-\i) at ({\s*\dx},{-\i*\dy}) {};}}

  % stage s pairs i with i XOR 2^(2-s); pgf has no xor, so the partners are
  % listed per stage -- which is exactly what the picture is about
  \foreach \a/\b in {0/4, 1/5, 2/6, 3/7}{
    \draw[wire] (p0-\a) -- (p1-\b);  \draw[wire] (p0-\b) -- (p1-\a);
    \draw[wire] (p0-\a) -- (p1-\a);  \draw[wire] (p0-\b) -- (p1-\b);}
  \foreach \a/\b in {0/2, 1/3, 4/6, 5/7}{
    \draw[wire] (p1-\a) -- (p2-\b);  \draw[wire] (p1-\b) -- (p2-\a);
    \draw[wire] (p1-\a) -- (p2-\a);  \draw[wire] (p1-\b) -- (p2-\b);}
  \foreach \a/\b in {0/1, 2/3, 4/5, 6/7}{
    \draw[wire] (p2-\a) -- (p3-\b);  \draw[wire] (p2-\b) -- (p3-\a);
    \draw[wire] (p2-\a) -- (p3-\a);  \draw[wire] (p2-\b) -- (p3-\b);}

  % redraw the points so the wires pass underneath
  \foreach \s in {0,...,3}{\foreach \i in {0,...,7}{
    \node[pt] at ({\s*\dx},{-\i*\dy}) {};}}

\end{tikzpicture}
\end{document}
